\documentclass[11pt]{article}
\usepackage[a4paper,margin=24mm]{geometry}
\usepackage{fontspec}
\usepackage{xcolor}
\usepackage{amsmath,amssymb}
\usepackage{booktabs,longtable,array}
\usepackage[hidelinks]{hyperref}
\setmainfont[Path=fonts/,Extension=.otf,UprightFont=*-Regular,BoldFont=*-Bold]{NotoSerifCJKkr}
\definecolor{reviewercolor}{RGB}{38,82,126}
\definecolor{responsecolor}{RGB}{36,113,78}
\definecolor{revisioncolor}{RGB}{145,82,24}
\setlength{\parindent}{0pt}
\setlength{\parskip}{0.55em}
\setlength{\fboxsep}{7pt}
\setlength{\emergencystretch}{3em}
\renewcommand{\familydefault}{\rmdefault}
\raggedright
\newcommand{\fieldheading}[2]{\par\medskip\noindent{\color{#1}\rule{3pt}{1.25em}}\hspace{0.55em}\textbf{#2}\par\smallskip}
\newenvironment{reviewitem}[1]{\subsection*{#1}}{\par\medskip\hrule\bigskip}
\newenvironment{reviewercomment}{\fieldheading{reviewercolor}{1. 리뷰어 코멘트}\begin{quote}}{\end{quote}}
\newenvironment{authorresponse}{\fieldheading{responsecolor}{2. 저자 답변}}{\par}
\newenvironment{manuscriptrevision}[1]{\fieldheading{revisioncolor}{3. 논문 반영 결과}\textbf{수정 위치:} #1\par\smallskip}{\par}
\newcommand{\manuscripttitle}{연속 인과관계(CoC) 주석의 시간·조건 일관성 검사 프레임워크: 적용 가능성과 근거 추적}
\newcommand{\manuscriptid}{확인 필요 (제공 기록에서 확인되지 않음)}
\newcommand{\revisionround}{확인 필요 (개정 소스 버전: revision-v01)}
\newcommand{\responsedate}{2026년 10월 12일}
% Generated from current manuscript auxiliary labels.
\expandafter\def\csname mloc@sec:intro\endcsname{1절(1쪽)}
\expandafter\def\csname mloc@fig:overview-pipeline\endcsname{그림 1(2쪽)}
\expandafter\def\csname mloc@sec:related\endcsname{2절(2쪽)}
\expandafter\def\csname mloc@sec:related-reasoning\endcsname{2.1절(2쪽)}
\expandafter\def\csname mloc@sec:related-rules\endcsname{2.2절(2쪽)}
\expandafter\def\csname mloc@sec:related-nl\endcsname{2.3절(2쪽)}
\expandafter\def\csname mloc@sec:related-rv\endcsname{2.4절(3쪽)}
\expandafter\def\csname mloc@sec:related-position\endcsname{2.5절(3쪽)}
\expandafter\def\csname mloc@sec:data\endcsname{3절(3쪽)}
\expandafter\def\csname mloc@sec:data-overview\endcsname{3.1절(3쪽)}
\expandafter\def\csname mloc@sec:provenance\endcsname{3.2절(3쪽)}
\expandafter\def\csname mloc@sec:review-protocol\endcsname{3.3절(3쪽)}
\expandafter\def\csname mloc@sec:problem\endcsname{4절(3쪽)}
\expandafter\def\csname mloc@tab:related-position\endcsname{표 1(4쪽)}
\expandafter\def\csname mloc@tab:evaluation-data\endcsname{표 2(4쪽)}
\expandafter\def\csname mloc@tab:provenance\endcsname{표 3(4쪽)}
\expandafter\def\csname mloc@sec:notation\endcsname{4.1절(4쪽)}
\expandafter\def\csname mloc@sec:error-types\endcsname{4.2절(4쪽)}
\expandafter\def\csname mloc@eq:hold-conflict\endcsname{식 2(4쪽)}
\expandafter\def\csname mloc@eq:missing-evidence\endcsname{식 6(4쪽)}
\expandafter\def\csname mloc@tab:model-scope\endcsname{표 4(5쪽)}
\expandafter\def\csname mloc@sec:semantics\endcsname{4.3절(5쪽)}
\expandafter\def\csname mloc@eq:active-update\endcsname{식 9(5쪽)}
\expandafter\def\csname mloc@sec:boundary-semantics\endcsname{4.4절(5쪽)}
\expandafter\def\csname mloc@tab:boundary-policy\endcsname{표 5(5쪽)}
\expandafter\def\csname mloc@sec:issue-definition\endcsname{4.5절(5쪽)}
\expandafter\def\csname mloc@tab:parser-rules\endcsname{표 6(6쪽)}
\expandafter\def\csname mloc@eq:response-contract\endcsname{식 13(6쪽)}
\expandafter\def\csname mloc@sec:method\endcsname{5절(6쪽)}
\expandafter\def\csname mloc@sec:workflow\endcsname{5.1절(6쪽)}
\expandafter\def\csname mloc@sec:extraction\endcsname{5.2절(6쪽)}
\expandafter\def\csname mloc@sec:models\endcsname{5.3절(6쪽)}
\expandafter\def\csname mloc@sec:queries\endcsname{5.3절(6쪽)}
\expandafter\def\csname mloc@sec:trace\endcsname{5.3절(6쪽)}
\expandafter\def\csname mloc@lst:condition-step\endcsname{목록 1(6쪽)}
\expandafter\def\csname mloc@fig:method-pipeline\endcsname{그림 2(7쪽)}
\expandafter\def\csname mloc@tab:worked-example\endcsname{표 7(7쪽)}
\expandafter\def\csname mloc@tab:model-variables\endcsname{표 8(7쪽)}
\expandafter\def\csname mloc@fig:condition-observer\endcsname{그림 3(7쪽)}
\expandafter\def\csname mloc@lst:queries\endcsname{목록 2(7쪽)}
\expandafter\def\csname mloc@tab:query-meaning\endcsname{표 9(7쪽)}
\expandafter\def\csname mloc@sec:boundary-implementation\endcsname{5.4절(7쪽)}
\expandafter\def\csname mloc@sec:model-scope\endcsname{5.5절(8쪽)}
\expandafter\def\csname mloc@sec:experiments\endcsname{6절(8쪽)}
\expandafter\def\csname mloc@sec:legacy-test\endcsname{6.1절(8쪽)}
\expandafter\def\csname mloc@sec:scenario-design\endcsname{6.2절(8쪽)}
\expandafter\def\csname mloc@sec:retest-design\endcsname{6.3절(8쪽)}
\expandafter\def\csname mloc@sec:boundary-design\endcsname{6.4절(8쪽)}
\expandafter\def\csname mloc@sec:bounded-design\endcsname{6.5절(8쪽)}
\expandafter\def\csname mloc@tab:comparison-paths\endcsname{표 10(9쪽)}
\expandafter\def\csname mloc@tab:scenario-variants\endcsname{표 11(9쪽)}
\expandafter\def\csname mloc@sec:human-design\endcsname{6.6절(9쪽)}
\expandafter\def\csname mloc@sec:aux-design\endcsname{6.7절(9쪽)}
\expandafter\def\csname mloc@tab:regression\endcsname{표 12(9쪽)}
\expandafter\def\csname mloc@sec:results\endcsname{7절(9쪽)}
\expandafter\def\csname mloc@sec:regression-results\endcsname{7.1절(9쪽)}
\expandafter\def\csname mloc@sec:first-results\endcsname{7.2절(9쪽)}
\expandafter\def\csname mloc@tab:first-results\endcsname{표 13(10쪽)}
\expandafter\def\csname mloc@tab:failures\endcsname{표 14(10쪽)}
\expandafter\def\csname mloc@sec:retest-results\endcsname{7.3절(10쪽)}
\expandafter\def\csname mloc@sec:agreement-results\endcsname{7.4절(10쪽)}
\expandafter\def\csname mloc@tab:retest\endcsname{표 15(10쪽)}
\expandafter\def\csname mloc@tab:implementation\endcsname{표 16(10쪽)}
\expandafter\def\csname mloc@tab:boundary-results\endcsname{표 17(10쪽)}
\expandafter\def\csname mloc@tab:bounded-results\endcsname{표 18(10쪽)}
\expandafter\def\csname mloc@sec:boundary-results\endcsname{7.5절(10쪽)}
\expandafter\def\csname mloc@sec:bounded-results\endcsname{7.6절(10쪽)}
\expandafter\def\csname mloc@tab:agreement\endcsname{표 19(11쪽)}
\expandafter\def\csname mloc@tab:reviewer-distribution\endcsname{표 20(11쪽)}
\expandafter\def\csname mloc@sec:human-results\endcsname{7.7절(11쪽)}
\expandafter\def\csname mloc@sec:trajectory-results\endcsname{7.8절(11쪽)}
\expandafter\def\csname mloc@sec:discussion\endcsname{8절(11쪽)}
\expandafter\def\csname mloc@sec:discussion-results\endcsname{8.1절(11쪽)}
\expandafter\def\csname mloc@sec:review-use\endcsname{8.2절(11쪽)}
\expandafter\def\csname mloc@tab:disagreement-examples\endcsname{표 21(12쪽)}
\expandafter\def\csname mloc@tab:trajectory-cases\endcsname{표 22(12쪽)}
\expandafter\def\csname mloc@tab:review-actions\endcsname{표 23(12쪽)}
\expandafter\def\csname mloc@sec:threats\endcsname{8.3절(12쪽)}
\expandafter\def\csname mloc@sec:future-work\endcsname{8.4절(12쪽)}
\expandafter\def\csname mloc@sec:conclusion\endcsname{9절(12쪽)}
\expandafter\def\csname mloc@sec:appendix-trajectory\endcsname{9절(12쪽)}
\expandafter\def\csname mloc@tab:limitations\endcsname{표 24(13쪽)}
\expandafter\def\csname mloc@tab:trajectory-details\endcsname{표 25(13쪽)}

% Generated from exact excerpts in the current main.pdf.
\expandafter\def\csname quotepages@R1-m3\endcsname{12쪽}
\expandafter\def\csname quotepages@R1-m2\endcsname{3쪽}
\expandafter\def\csname quotepages@R1-m1\endcsname{3쪽}
\expandafter\def\csname quotepages@R1-M6\endcsname{3쪽}
\expandafter\def\csname quotepages@R1-M5\endcsname{11쪽}
\expandafter\def\csname quotepages@R1-M4\endcsname{10쪽}
\expandafter\def\csname quotepages@R1-M3\endcsname{11쪽}
\expandafter\def\csname quotepages@R1-M2\endcsname{6쪽}
\expandafter\def\csname quotepages@R1-M1\endcsname{12쪽}
\expandafter\def\csname quotepages@R2-7\endcsname{2쪽}
\expandafter\def\csname quotepages@R2-6\endcsname{3쪽}
\expandafter\def\csname quotepages@R2-5\endcsname{1쪽}
\expandafter\def\csname quotepages@R2-4\endcsname{2쪽}
\expandafter\def\csname quotepages@R2-3\endcsname{1쪽}
\expandafter\def\csname quotepages@R2-2\endcsname{1쪽}
\expandafter\def\csname quotepages@R2-1.2\endcsname{1쪽}
\expandafter\def\csname quotepages@R2-1.1\endcsname{1쪽}
\expandafter\def\csname quotepages@R3-11\endcsname{11쪽}
\expandafter\def\csname quotepages@R3-10\endcsname{9쪽}
\expandafter\def\csname quotepages@R3-9\endcsname{11쪽}
\expandafter\def\csname quotepages@R3-8\endcsname{10쪽}
\expandafter\def\csname quotepages@R3-7\endcsname{10쪽}
\expandafter\def\csname quotepages@R3-6\endcsname{6쪽}
\expandafter\def\csname quotepages@R3-5\endcsname{6쪽}
\expandafter\def\csname quotepages@R3-4\endcsname{5쪽}
\expandafter\def\csname quotepages@R3-3\endcsname{4쪽}
\expandafter\def\csname quotepages@R3-2\endcsname{3쪽}
\expandafter\def\csname quotepages@R3-1\endcsname{3쪽}

\newcommand{\mloc}[1]{\csname mloc@#1\endcsname}
\newcommand{\quotepages}[1]{\csname quotepages@#1\endcsname}
\newcommand{\responsefront}[2]{%
\begin{center}{\LARGE\bfseries 심사 의견 답변서\par}\vspace{1em}{\large\bfseries\manuscripttitle\par}\end{center}
\vspace{0.8em}
\begin{tabular}{@{}ll@{}}논문 접수번호: & \manuscriptid \\ 수정 차수: & \revisionround \\ 답변일: & \responsedate\end{tabular}
\section*{답변에 앞서}
편집위원장님과 심사위원님께,

연속 CoC의 검사 범위, 의미 변환과 평가 근거를 구체화할 수 있도록 의견을 주셔서 감사합니다.
이번 개정은 자연 오류 성능 대신 시간·조건 관계의 검사 적용 가능성과 근거 추적으로 주장을 제한했습니다.
기존 통제 시나리오의 최초 결과와 개발 후 동일 사례 재시험을 구분하고, 81건의 장면 경계 시험과 기존 사람 라벨의 AC1 재분석을 본문에 추가했습니다.
선별 없는 유한 직적 전수 시험은 고정 XML의 step 코드와 별도 수식 판정기를 대조한 구현 적합성 근거입니다.
두 경로에 같은 작성자가 관여했으므로 외부 독립 시험이나 자연 의미 정답 검증을 완료했다고 주장하지 않습니다. 사람 재평가도 수행하지 않았습니다.
개정 원고의 파란색은 원 저장소의 최초 제출본 대비 이번에 추가·수정한 내용을 나타냅니다.
검정색 clean 원고는 동일 내용입니다. 아래 절·쪽·표 번호는 이번 개정 원고를 기준으로 합니다.
접수번호와 실제 심사 수정 차수는 기록이 없어 확인 필요로 남겼습니다.

\section*{심사위원 #1}
\textbf{총평(요지)}\par
\begin{quote}#2\end{quote}}


\begin{document}

\responsefront{1}{연속 요구의 상태 관리와 최초 위반 추적이라는 연구 방향은 의미가 있으나, 기대 판정을 확인한 합성 40쌍과 낮은 사람 일치도만으로 자연 오류 성능을 입증하기 어렵다. 독립 검증, 규칙 추출, 검토 신뢰도, 비교기와 재현성 보완이 필요하다.}

\clearpage

\begin{reviewitem}{R1 Major 1. 자연 오류 검증과 주장 범위}

\begin{reviewercomment}

현재 핵심 결과인 prior-requirement checker 40/40, current-event-only checker 20/40은 의미 있는 ablation이지만, 40개 모두 기대 결과가 나타나도록 구성된 합성 사례입니다. 따라서 실제 자연 데이터에서 제안 방법이 얼마나 잘 동작하는지 평가할 수 없기 때문에, 가능하다면 무작위 또는 사전에 정의한 sampling protocol을 사용하여 자연 발생 consistent/inconsistent CoC ground-truth set을 구축하고 TP, FP, TN, FN, sensitivity/recall, specificity, precision, F1-score 및 95\% confidence interval 등을 제시할 필요가 있습니다. 자연 오류 자료 구축이 현실적으로 어렵다면 제목, 초록, 결론의 claim을 ``rule execution/feasibility framework'' 수준으로 명확히 제한해야 합니다.

\end{reviewercomment}

\begin{authorresponse}

지적을 수용하고 심사위원님이 제시한 적용 가능성(feasibility) 수준의 대안을 선택했습니다. 이번 작업에서 자연 오류 정답을 새로 수집하지 않았으므로 자연 sensitivity, specificity, precision/recall/F1과 신뢰구간을 만들지 않았습니다. 제목에 적용 가능성과 근거 추적을 명시하고 초록·서론·결론을 같은 범위로 수정했습니다.

통제 자료는 설계자가 작성한 144개 조합입니다. 기존 모델의 최초 판정 일치는 96/144이고, 실패 48개를 보존했습니다. 실패 분석 후 조건 연결 모델의 144/144는 같은 사례에 대한 개발 후 재시험입니다. 독립 자연 오류 성능으로 제시하지 않습니다. 단순 FSM도 144/144이므로 UPPAAL 정확도 우위도 주장하지 않습니다.

% Supplemental evidence 2026-10-12
별도 81건의 장면 경계 시험도 본문에 통합했습니다. 유지 모델은 81/81, 초기화는 14/81입니다. 초기화 모델은 확정 모순 41/44와 근거 부족 26/29를 양립으로 판정했습니다. 두 결과는 작성 시험입니다. 선별 없는 유한 범위의 35,042,928건을 별도 수식 계산과 고정 XML의 상태 갱신 코드로 비교한 결과 불일치 0이었지만, 이는 명세 적합성의 부분 근거이며 자연 오류 성능이나 외부 독립 의미 정답이 아닙니다.
% End supplemental evidence

\end{authorresponse}

\begin{manuscriptrevision}{\mloc{sec:conclusion}; 발췌: \quotepages{R1-M1}}
본문의 실제 수정 문장은 다음과 같습니다.

\begin{quote}
현재 결과는 명시된 프레임의 시간ㆍ조건 관계를 실행하고 검토 근거를 추적하는 방법 실증이다.
\end{quote}
\end{manuscriptrevision}

\end{reviewitem}

\clearpage

\begin{reviewitem}{R1 Major 2. 원문에서 검사 입력으로의 변환}

\begin{reviewercomment}

프레임워크의 핵심은 자연어 CoC에서 action, continuing requirement, end condition, action order, True/False/Unknown evidence를 추출하는 과정인데, 이 변환이 자동인지 수동인지, rule-based parser인지 LLM 기반인지, 동의 표현과 부정문·조건문·복합행동·복수 requirement를 어떻게 처리하는지, end condition이 불명확할 때 Unknown으로 분류하는 기준이 무엇인지 충분히 설명되지 않고 있습니다. 실제 CoC 원문 → structured representation → checker input으로 이어지는 사례를 표 또는 pseudocode 형태로 제시하는 것이 필요합니다.

\end{reviewercomment}

\begin{authorresponse}

기존 자동 경로를 정규식 기반의 제한된 영어 파서로 명시했습니다. 행동 동의어 묶음, 명령 문맥, then/and then 단계 분리, until 해제 조건, once/after 계기, 대상 누락과 순서 모호 처리 규칙을 표로 정리했습니다. 조건의 참·거짓 증거는 자동 확정하지 않으며, 자동 텍스트 경로의 144건 전부 미상 출력 결과를 5.2절 첫머리에도 밝혔습니다.

144건의 주평가 입력은 설계자가 원문과 함께 작성한 프레임입니다. 작성 보행자 사례에서 공통 e0과 세 e1 변형의 원문·의무 A·조건 A/B·참/거짓/미상·사례 전체 관측 배열·다섯 질의·판정을 한 표로 연결했습니다. 프레임의 참·거짓은 작성 문장이 명시한 상태이지 독립 사람 정답이 아닙니다. 사건 고정, 프레임 구성, 판정과 기록의 세 단계를 서론 및 그림과 맞추고, 최초 문제의 사건·의무에서 원문 구절과 연결하여 제시하는 과정을 설명했습니다.

% Supplemental evidence 2026-10-12
자연 CoC 90장면의 원문 399건에 대한 파서 집계를 추가했습니다. 행동 인식 원문은 230건, 미인식은 169건이며 265단계 중 구문 해석 완료 219·대상 누락 44·순서 모호 2건입니다. until 해제 조건은 0건, once/after 계기는 10건입니다. scene-0051의 실제 문장은 자동 처리 사례로 별도 제시했습니다. 이는 구문 처리 집계이지 의미 추출 정확도가 아닙니다. 별도 81건은 원 자연어 구절이 없는 기호 입력이며, 대상 동일성·적용 범위·증거 유효성은 설계자 제공 값입니다.
% End supplemental evidence

\end{authorresponse}

\begin{manuscriptrevision}{\mloc{sec:extraction}; 발췌: \quotepages{R1-M2}}
본문의 실제 수정 문장은 다음과 같습니다.

\begin{quote}
기존 자동 변환은 LLM 호출이나 수작업 정답 입력을 쓰지 않는 정규식 기반 영어 파서다. 이 파서를 거친 기존 자동 텍스트 경로는 144건 모두 미상 출력(\texttt{UNKNOWN})을 반환했다(7.2절).
\end{quote}
\end{manuscriptrevision}

\end{reviewitem}

\clearpage

\begin{reviewitem}{R1 Major 3. 검토자 일치도와 합의 기록}

\begin{reviewercomment}

텍스트 일관성 Fleiss' kappa=0.027, temporal requirement kappa=0.203, required action sequence kappa=0.291은 평가 기준의 안정성과 재현성을 평가하는 데 적절한지 의문입니다. 검토자의 전문성, 사전 교육 여부, annotation guideline, reviewer별 판정 분포, disagreement 유형, consensus 형성 절차를 상세히 제시하는 것이 필요합니다. 또한, 가능하다면 guideline 수정 후 재평가하거나 독립 adjudicator를 활용하여 평가 기준의 명확성을 보강할 수 있습니다.

\end{reviewercomment}

\begin{authorresponse}

낮은 일치도가 자연 오류 정답의 신뢰도를 제한한다는 지적에 동의합니다. 네 원 CSV를 다시 집계하여 텍스트 κ=0.027, 시간 요구 κ=0.203, 행동열 κ=0.291과 항목별 쌍별 일치율·불일치 장면 수를 제시했습니다. 검토자별 텍스트 판정 분포도 추가했습니다. 텍스트 불일치는 23/27이고 모순 판정 수는 검토자별 10/0/1/1건입니다.

지침과 원 메모에서 확인되는 판단 차이는 기술했으나 원인별 확정 건수는 만들지 않았습니다. 전문성·교육·독립 중재·당시 논의 절차는 충분히 복원되지 않았고 합의 CSV의 사유 메모는 0/27입니다. 새 사람 재평가나 중재를 수행하지 않았습니다. 합의 결과를 객관적 자연 정답으로 사용하지 않으며 신뢰도 개선 요구는 남은 한계로 인정합니다.

보존된 검토 지침과 A/B HTML에서 원문·상대 시각만 제공하고 후보/비교 구분·원 장면 ID·검사기 결과를 표시하지 않은 것을 확인했습니다. 집단 대응표는 별도 파일이고, A/B는 같은 27건을 서로 다른 고정 순서로 제시합니다. 제출본의 ``모른 채''는 확인 가능한 범위인 ``1차 검토 화면에서 집단 구분을 표시하지 않았다''로 구체화했습니다.

% Supplemental evidence 2026-10-12
원 4인×27장면의 AC1을 추가 재분석했습니다. 시간 요구 0.627, 행동열 0.323, 텍스트 0.301이며 기존 κ를 대체하지 않습니다. 시간 요구는 YES 81/108의 범주 편중으로 우연 일치 항에 민감합니다. 텍스트의 원시 일치 0.407과 불일치 23/27은 그대로이고, 행동열은 39관측 문자열 범주에 조건부입니다. 새 사람 평가·지침 개선 효과·자격/교육/합의 기록 복원을 주장하지 않습니다.
% End supplemental evidence

\end{authorresponse}

\begin{manuscriptrevision}{\mloc{sec:human-results}; 발췌: \quotepages{R1-M3}}
본문의 실제 수정 문장은 다음과 같습니다.

\begin{quote}
AC1 추가 재분석은 세 항목 모두에서 Fleiss 계수보다 높았다. 범주 편중에 따른 보정 민감도이며 지침의 안정성이 개선됐다는 뜻이 아니다.
\end{quote}
\end{manuscriptrevision}

\end{reviewitem}

\clearpage

\begin{reviewitem}{R1 Major 4. 비교 방법의 확장}

\begin{reviewercomment}

현재 비교 대상인 current-event-only checker는 구조상 cross-event error를 검출할 수 없으므로 제안 방법의 필요성을 보여 주는 ablation으로는 적절하지만, 충분한 경쟁 baseline이라고 보기는 어렵습니다. fixed previous-N-event context, simple finite-state-machine, sequence-level rule-based checker, temporal logic/LTL 기반 방법 또는 full sequence를 입력받는 LLM consistency judge 등과의 비교를 고려할 필요가 있습니다.

\end{reviewercomment}

\begin{authorresponse}

현재 사건 검사를 기억 제거 비교(ablation)로 한정하고 이전 1/3/5개 사건 창과 별도 단순 FSM을 추가한 완료 실험을 반영했습니다. 144건의 최초 일치는 현재 사건 72, 이전 창 84/90/90, 기존 전체 이력 96, FSM 144개입니다. 간격별 결과와 최초 위반 위치도 제시했습니다.

창 비교는 같은 기존 검사기에 잘린 이력을 제공한 것이고, FSM은 조건별 프레임을 입력받습니다. 기존 경로의 단일 의무 투영 손실과 미상 누적 정책이 달라 엄밀한 정확도 우위 비교로 쓰지 않았습니다. FSM도 전 사례에 성공했으므로 UPPAAL만이 필요한 방법이라고 하지 않습니다. 범용 LTL 비교나 전체 문장열 LLM 판정기는 이번 실험에 포함하지 않았습니다.

LLM 판정기 비교의 미실시 사실을 평가 설계 \mloc{sec:scenario-design}에도 명시했습니다. \mloc{sec:threats}에는 이 비교가 없어 범용 자연어 일관성 판정기에 대한 상대 성능을 판단할 수 없음을 추가했습니다.

% Supplemental evidence 2026-10-12
81건의 유지/초기화 비교는 경계 기억의 제거 비교로 추가했습니다. FSM과 이전 N개 사건 창의 144건 비교를 대체하지 않으며, 다른 자료의 수치를 합쳐 정확도 우위를 계산하지 않습니다. 전체 문장열 LLM 판정기는 후속 작업으로 남아 있습니다.
% End supplemental evidence

\end{authorresponse}

\begin{manuscriptrevision}{\mloc{sec:first-results}; 발췌: \quotepages{R1-M4}}
본문의 실제 수정 문장은 다음과 같습니다.

\begin{quote}
FSM은 144/144, 최초 위반 72/72였다. UPPAAL이 단순 비교기보다 정확하다는 근거는 없다.
\end{quote}
\end{manuscriptrevision}

\end{reviewitem}

\clearpage

\begin{reviewitem}{R1 Major 5. 반응 시간 설정과 탐색 규칙}

\begin{reviewercomment}

반응시간 분석의 W=0.5 s, delta=0.1 m/s, rho=0.7, D=3 s 및 STOP/YIELD에서 v≤0.5 m/s 기준은 외부 안전 기준이나 법규에서 직접 도출된 값이 아닙니다. 특히, 일부 규칙은 초기 후보 사례 검토 이후 추가된 탐색 규칙이므로 exploratory analysis임을 명확히 해야 합니다. threshold 선택의 외부 근거 또는 별도의 sensitivity analysis를 보강하고, 130개 중 58개 사건이 27개 설정에 따라 결과가 달라졌다는 점을 명확히 기술해야 합니다.

\end{reviewercomment}

\begin{authorresponse}

외부 안전 기준에서 도출한 설정이 아님을 수용했습니다. 속도 기반 반응 검사를 탐색적 보조 분석으로 명시하고 W, δ, ρ, D와 이미 정지 규칙의 적용 순서를 설명했습니다. 이미 정지 규칙은 초기 후보를 본 뒤 추가한 사실도 유지했습니다.

평가 설계 \mloc{sec:aux-design}에 민감도 분석의 조합과 적용 대상을 명시했습니다. 탐지 대상 130건에는 $D=3$초로 고정하고 $W\in\{0.25,0.5,1.0\}$초, $\delta\in\{0.05,0.1,0.2\}$~m/s, $\rho\in\{0.5,0.7,0.9\}$의 27개 조합을 적용했습니다. $D=1,2,3,4,5$초를 포함한 사건당 135개 조합은 기본 설정에서 제한 시간 안에 반응을 확인하지 못한 5건에만 적용한 추가 분석이며 130건 전체의 제한 시간 분석으로 제시하지 않습니다.

기존 27설정 결과를 재분석하면 130건은 항상 검출 72, 설정 의존 57, 항상 미검출 1입니다. 따라서 심사 당시 원고의 58건 전체를 설정 의존이라고 쓴 표현을 정정했습니다. 기본 분모 134=125+4+5와 탐지 대상 130의 차이도 설명했습니다. 이 민감도 결과를 법규 근거나 물리적 안전 기준으로 해석하지 않습니다.

\end{authorresponse}

\begin{manuscriptrevision}{\mloc{sec:trajectory-results}; 발췌: \quotepages{R1-M5}}
본문의 실제 수정 문장은 다음과 같습니다.

\begin{quote}
탐지 대상 130건을 27설정에서 재집계하면 항상 검출 72, 설정 의존 57, 항상 미검출 1이다. 확인되지 않은 5건은 제한 시간을 4초 이상으로 늘리면 반응이 검출되는 조합이 크게 늘었으나(5사건$\times$27설정 중 $D=3$초 18개, $D=4$초 80개, $D=5$초 103개), 한 장면은 5초에서도 기본 설정으로 검출되지 않았다.
\end{quote}
\end{manuscriptrevision}

\end{reviewitem}

\clearpage

\begin{reviewitem}{R1 Major 6. 재현성과 자료 출처}

\begin{reviewercomment}

CoC 생성 모델/version, prompt 또는 configuration, quality-score threshold, 403→290 선별 기준, Python checker source code, UPPAAL XML, synthetic-case generator, 40 reference-modified pairs, annotation guideline 및 27개 human review raw label 등을 가능한 범위에서 공개할 필요가 있습니다. 특히 논문에서 일부 생성 모델과 자동 표시 도구의 정확한 버전이 남아 있지 않다고 기술한 점은 재현성 측면에서 매우 중요한 한계점입니다.

\end{reviewercomment}

\begin{authorresponse}

복원 사실과 미기록을 분리했습니다. 자료 설명서의 Qwen/Qwen3.5-35B-A3B 계열은 확인되지만 생성 실행별 모델 버전, 원 생성 지시문·설정과 품질 점수의 계산식·통과 기준은 복원되지 않았습니다. 403→290은 저장된 품질 통과 표시를 사용한 선정이고, 290→134는 방향이 불명확하거나 복합인 156건을 제외한 절차입니다.

3장 첫머리에 전체 평가 자료 개요를 두고 자연 자료의 출처·선정과 사람 검토 기록을 분리했습니다. 재현용 시나리오·코드·UPPAAL 모델·질의·실행 기록의 보존 범위와 원 CoC·개별 검토 CSV의 접근 제한을 명시했습니다. 이번 재현 패키지의 외부 공개 범위·배포 주소는 미확정이며, 저장소의 공개용 분류를 실제 배포 완료로 제시하지 않습니다. 완전한 원 자료 재생성을 주장하지 않습니다.

% Supplemental evidence 2026-10-12
신규 경계 시험과 유한 범위 적합성 시험에는 입력 프로토콜, 고정 소스·XML·판정기 해시, 전체 생성 규칙과 불일치를 보존했습니다. 적합성 시험은 실제 XML에서 추출한 상태 갱신 코드를 C++로 직접 실행한 것으로서 UPPAAL 탐색기를 새로 실행했다는 뜻은 아닙니다. 과거 생성 모델·prompt/config와 검토자 기록의 결손은 그대로 인정합니다.
% End supplemental evidence

\end{authorresponse}

\begin{manuscriptrevision}{\mloc{sec:provenance}; 발췌: \quotepages{R1-M6}}
본문의 실제 수정 문장은 다음과 같습니다.

\begin{quote}
CoC-Nusc의 배포 버전, 생성 실행별 모델 버전, 생성 지시문과 설정은 복원하지 못했다.
\end{quote}
\end{manuscriptrevision}

\end{reviewitem}

\clearpage

\begin{reviewitem}{R1 Minor 1. 자연 후보 15장면의 선정}

\begin{reviewercomment}

27개 자연 검토 장면 중, 15개 automatically selected candidate scenes의 선정 기준을 명확히 기술해야 합니다. 비무작위 표본이라는 사실뿐 아니라 candidate selection algorithm 또는 조건 자체가 결과에 영향을 줄 수 있습니다.

\end{reviewercomment}

\begin{authorresponse}

보존된 원 선정 알고리즘을 본문에 명시했습니다. 첫 문장의 정지·양보 표현과 둘째 문장의 가속·진행 표현이 있는 이웃쌍 23개를 찾고 15장면에서 가장 이른 창 하나씩을 선택합니다. 비교군은 후보 장면을 제외한 정지·양보 후 비진행 창으로, 사건 수·정지/양보 종류·원 위치 차이와 고정 시드 해시의 동률 처리로 대응시켜 중복 없는 12장면을 선택했습니다.

선정 규칙을 검토 결과에 맞춰 새로 만든 것이 아닙니다. 보존된 선정 프로그램의 재실행에서 원 27개 검토 ID와 장면·집단 구분이 재현됩니다. 이 표본은 특정 행동 전환을 더 포함하는 비무작위 표본이며 자연 모집단 오류율 추정에 사용하지 않습니다.

후보군·비교군 구분은 별도 대응표에만 있고 보존된 1차 검토 화면에는 표시되지 않습니다. 이 정보 가림을 3.3절에 명시했으며, 표본 선정의 비무작위성과 구분했습니다.


\end{authorresponse}

\begin{manuscriptrevision}{\mloc{sec:provenance}; 발췌: \quotepages{R1-m1}}
본문의 실제 수정 문장은 다음과 같습니다.

\begin{quote}
후보 선택기는 첫 문장에 정지ㆍ양보(\texttt{stop/yield}), 둘째 문장에 가속ㆍ진행(\texttt{accelerate/proceed}) 표현이 있는 쌍을 찾는다. 23쌍ㆍ15장면에서 장면별 가장 이른 창 하나를 선정했다.
\end{quote}
\end{manuscriptrevision}

\end{reviewitem}

\clearpage

\begin{reviewitem}{R1 Minor 2. 초록의 검토 분모}

\begin{reviewercomment}

초록의 ``309 adjacent pairs from 90 scenes, including 27 groups'' 표현은 309개 전체가 사람 검토된 것으로 오해될 수 있으므로, 수정이 필요합니다.

\end{reviewercomment}

\begin{authorresponse}

309인접쌍 전체를 사람이 검토한 것으로 읽히는 초록 표현을 제거하고, 3장의 자료 개요와 선정 설명에서 평가 단위를 구분했습니다. 자차 궤적과 대응되는 94장면·403사건 중 문장이 둘 이상인 90장면에서 309인접쌍을 만들었습니다. 이 쌍들은 후보 선정용이며 사람이 검토한 것은 별도로 선정한 27장면 묶음입니다. 309쌍 전체의 오류율을 보고하지 않습니다.


\end{authorresponse}

\begin{manuscriptrevision}{\mloc{sec:provenance}; 발췌: \quotepages{R1-m2}}
본문의 실제 수정 문장은 다음과 같습니다.

\begin{quote}
309는 후보 선정에 사용한 인접 문장 쌍의 수이고, 27은 사람이 검토한 장면 묶음의 수다. 여러 사건이 한 장면에 속하므로 사건ㆍ쌍ㆍ장면을 독립 시행으로 합산하지 않는다.
\end{quote}
\end{manuscriptrevision}

\end{reviewitem}

\clearpage

\begin{reviewitem}{R1 Minor 3. 결론의 필요성과 효과 표현}

\begin{reviewercomment}

결론의 이전 요구를 기억하는 검사가 필요하고 효과적이라는 표현은 현재 증거 수준보다 다소 넓게 읽힐 수 있으므로, 본 연구에서 정의한 cross-event synthetic error를 탐지하는 데 prior-requirement retention이 필요하였다 정도로 제한하는 것이 더 적절해 보입니다.

\end{reviewercomment}

\begin{authorresponse}

지적을 수용하여 필요성 표현을 작성된 미해제 충돌 사례에 한정했습니다. 결론에는 기존 검사기의 최초 실패(96/144)와 FSM의 동일 판정을 포함하고, 조건 연결 모델의 결과는 동일 사례 재시험이라고 다시 명시했습니다. 자동 텍스트 경로의 결과는 5.2절과 7.2절에 보고했습니다. 모든 CoC에 대한 필요성·효과, 자연 오류 성능, 학습 개선 또는 UPPAAL 정확도 우위로 확대하지 않습니다.

\end{authorresponse}

\begin{manuscriptrevision}{\mloc{sec:conclusion}; 발췌: \quotepages{R1-m3}}
본문의 실제 수정 문장은 다음과 같습니다.

\begin{quote}
작성된 미해제 충돌 사례에서는 이전 의무의 보존이 해당 사건 관계를 판정하는 데 필요했다. 다만 이 관찰을 모든 CoC에 대한 필요성ㆍ효과나 UPPAAL의 정확도 우위로 일반화하지 않는다.
\end{quote}
\end{manuscriptrevision}

\end{reviewitem}

\section*{맺음말}

심사위원님의 의견에 따라 주장 범위를 적용 가능성 실증으로 제한하고, 원문에서 검사 입력으로의 변환 경로, 비교 방법(이전 사건 창과 조건별 FSM), 반응 시간 설정의 민감도, 자료 출처와 선정 절차를 보완했습니다. 개발 후 동일 사례 재시험과 별도로 선별 없는 유한 범위 전수 대조를 추가했습니다. 다만 자연 오류에 대한 독립 검증, 검토자 일치도의 개선, 과거 생성 이력의 복원은 이번 개정에서 완료하지 못했으며, 이를 본문의 적용 범위와 후속 과제로 명시했습니다.

\end{document}
